\section{Analytical comparison and research boundary}\label{sec:scope}

\subsection{What the two constructions establish}
Both branches now have complete analytical realizations within their
chosen supplied-data models, together with internal assessments accepting
those scopes without required mathematical correction
\cite{TNDDirectReview,TNDAuxReview}. The common geometry and
source-transport comparison has its separate shared assessment
\cite{TNDSharedReview}. The table describes the mathematical status at
6 November 2026; it does not replace either branch's live research register.

\begin{table}[ht]
\small\centering
\begin{tabularx}{\textwidth}{@{}>{\raggedright\arraybackslash}p{0.19\textwidth}YY@{}}\toprule
Feature & Direct Newman realization & Auxiliary Hamiltonian dynamics \\\midrule
Prepared object & Normalized square root $f_{1/2}$ with complete angular density $R$. & Linear amplitude $F_\lambda B$ in a passive reference bank. \\
Evolution & Free unitary and the specified count instrument; condition on no count. & Free left heat $\ee^{-\tau L_{\rm left}}$ after each supplied preparation. \\
Newman parameter & The declared decreasing clock $\lambda(t)=1/2-\gamma t$ on $[-1/2,1/2]$. & An independent source parameter on any fixed finite compact interval. \\
Signed readout & Density cosine in $Q_t$, calibrated by joint success frequencies. & Phase-sensitive linear Fourier pairing with known absolute gain. \\
Exact result & $H_\lambda=H_{1/2}(0)(q_+-q_-)$. & $\Obs=q(\tau)\ee^{-\tau z^2}H_\lambda$. \\
Remaining input & Source loader, imprint, angular preparation, initial scale, quadrature access. & Source/imprint, matched passive bank, calibrated phase-sensitive law. \\
Current mathematical limit & No signed estimate follows from conditioning or positivity. & Zeros and collisions are preserved; inverse amplification grows as $\ee^{\tau k^2}$. \\
\bottomrule\end{tabularx}
\caption{Accepted conditional representations and their different burdens.}
\end{table}

The analytical progress is substantive at the level of domains and
representation: the square-root preparation has controlled derivatives,
all counts are retained, the strong law uses a stated initial and final
domain, and the auxiliary response has paid exponential moments and
actual zero-time derivatives. It also makes the remaining issue precise.
A positive source, positive angular state or positive operation does not
control a signed cosine cancellation. A new assertion about the original
zeros requires additional information about the actual source and
observation, beyond either exact representation.

A later source refresh during manuscript production found two branch
continuations. The direct branch completed initial collision reconnaissance
at proof-author/diagnostic scope; mechanism selection remains open. The
auxiliary branch completed a conceptual choice and issued an unactivated
execution request for first-derivative recovery conditioning. These later
items do not inherit the earlier realization assessments automatically.

\subsection{Direct continuation: count sectors and collision anatomy}
The subsequent direct note \cite{TNDDirectCollision} records the following
identities of the already specified instrument. They retain that note's
proof-author status. Put
\[
 w_t=F_{\lambda(t)}/Z_{\lambda(t)},\quad
 C(t,k)=\int_\R w_t(v)\cos(kv)\dd v,\quad
 d_0(t,k)=p_tC(t,k)=H_{\lambda(t)}(k)/H_{1/2}(0).
\]
Differentiating the complete source integrals gives
\begin{equation}\label{eq:diagnosticPDE}
 \partial_tw_t=-\gamma(v^2-\mu_2(t))w_t,\quad
 C_t=\gamma(C_{kk}+\mu_2C),\quad
 (d_0)_t=\gamma(d_0)_{kk},\qquad
 \mu_2(t)=\int_\R v^2w_t(v)\dd v.
\end{equation}
The physical-time sign is positive because $\lambda(t)$ decreases.
The normalized state and success mass remain regular even if the
signed readout has a multiple zero.

For each count output $\sigma_j(t)=K_j(t,0)\rho_*K_j(t,0)^*$ define
$d_j(t,k)=\Tr[\sigma_j(t)\cos(kQ_t)]$. Direct substitution of
\eqref{eq:Kraus} gives
\begin{equation}\label{eq:countjets}
 d_j=\frac{(-\gamma t)^j}{j!}\partial_k^{2j}d_0,\qquad
 |d_j|\le\Tr\sigma_j,\qquad
 \sum_{j\ge0}\eta^jd_j(t,k)
 =\frac{H_{\lambda(t)+\eta\gamma t}(k)}{H_{1/2}(0)},\quad0\le\eta\le1.
\end{equation}
Absolute convergence follows from $\sum_j\Tr\sigma_j=1$; inside the
source integral it is the positive exponential series. Thus
$\sum_jd_j(t,k)=C(0,k)$, and for $t>0$,
\begin{equation}\label{eq:counttransfer}
 \partial_td_j=\frac jt d_j-\frac{j+1}{t}d_{j+1},\qquad
 \partial_td_0=-d_1/t.
\end{equation}
These are analytical traces of existing outputs. The no-count readout
contract does not by itself grant experimental access to the same
observation after every nonzero count. Also, conservation of this
count-averaged cosine trace does not make the entire unconditioned
state stationary.

At a hypothetical interior nondegenerate double event, set
$a_*=(d_0)_{kk}(t_*,k_*)\ne0$. Taylor expansion using
\eqref{eq:diagnosticPDE} yields
\begin{equation}\label{eq:localdouble}
 d_0(t_*+h,k_*+x)
 =\frac{a_*}{2}(x^2+2\gamma h)
   +O(|x|^3+|hx|+h^2).
\end{equation}
The constants are local to the event. The real roots on the $h<0$ side
have leading displacement $\pm\sqrt{-2\gamma h}$; on the $h>0$ side
the corresponding complex displacements are
$\pm\ii\sqrt{2\gamma h}$. This follows by parabolic rescaling and
the analytic implicit-function theorem at the two simple scaled roots.
At the event $q_+=q_-=p_{t_*}/2>0$, while
$d_1(t_*,k_*)=-\gamma t_*a_*\ne0$. The singular behavior is in zero
tracking, not in the positive state or success mass. This conditional
local description neither asserts a theta collision nor covers higher
multiplicities by a generic-double assumption.

The same direct note checks a reflected super-Gaussian collision
control inside the physical slab, at $\lambda=1/4$, $k=\pi/4$, using
its own smooth-square-root tail bounds. It satisfies these generic
balance laws while failing the canonical theta arithmetic expansion.
Accordingly the count-sector identities organize the same source
moments but supply no separating signed bound. The next direct
conceptual step remains the search for a source-specific mechanism;
no follow-up estimate is fixed by this manuscript.

\subsection{Auxiliary continuation: the value-and-derivative recovery map}
The auxiliary decision \cite{TNDAuxConditioning} fixes a finite-window
conditioning diagnostic. To display its background algebra, take a
fixed derivative scale $\sigma>0$ and define
\[
 \mathbf h=\binom{H_\lambda(k)}{\sigma H'_\lambda(k)},\qquad
 \mathbf o=\binom{\Obs_{\lambda,\tau}(k)}{\sigma\partial_k\Obs_{\lambda,\tau}(k)},
 \qquad m_\tau(k)=q(\tau)\ee^{-\tau k^2}.
\]
One differentiation of the already accepted response gives
\begin{equation}\label{eq:jetmap}
 \mathbf o=m_\tau(k)
 \begin{pmatrix}1&0\\-2\sigma\tau k&1\end{pmatrix}\mathbf h.
\end{equation}
Thus recovery must account for derivative mixing as well as angular and
Gaussian attenuation. The issued execution request asks for the sharp
Euclidean inverse norm on a fixed finite real window and a comparison
with direct recovery under a declared additive pair-error hypothesis.
Those sharp bounds remain proof targets at this source capture; they
are not asserted here as completed research.

Covariant derivative processing removes the signal's triangular factor,
 but transforms an isotropic raw error ball into an ellipsoid. For an
uncertainty set transported by the exact invertible map, applying its
inverse returns the original set. This algebra supplies no gain in
information by itself. The selected diagnostic introduces neither a
physical noise law nor a positive canonical value-and-derivative margin.
It does not change the supplied apparatus or prove a collision restriction.

\subsection{Controls and the scope of the conclusions}
The repository's collision-control collection \cite{CollisionControls}
provides a premise screen for any later sign claim. The generic Fourier,
Green, count-filter and heat identities can hold for other sources that
admit collisions. Their applicability must be checked against the
particular source regularity and preparation interval; the table does
not assert that every control meets every domain assumption.

\begin{center}\small
\begin{tabularx}{\textwidth}{@{}>{\raggedright\arraybackslash}p{0.25\textwidth}Y@{}}\toprule
Control class & Premise absent from the canonical application; surviving structure \\\midrule
Euler-factor perturbations & The unchanged complete theta arithmetic source is absent. Applicable smooth positive even sources still admit generic Fourier/Green constructions and conditional filtering. \\
Weighted-square and finite-prime perturbations & Complete unit coefficients and full canonical reflection fail. Positive half-line sums and finitely matched endpoints do not establish a smooth full-line square root. \\
Reflected super-Gaussian and smooth positive-primitive controls & The exact theta arithmetic expansion fails. Where the required tails and domains hold, smoothness, positivity, reflection and the representation identities survive. \\
Atomic or nonsmooth finite-profile controls & Fourier observations and finite moments may exist while the smooth source and graph-domain hypotheses fail or require a separate construction. \\
\bottomrule\end{tabularx}
\end{center}

This screen is analytical reuse of the recorded controls, not a new
control computation. No control is required to violate a generic
representation theorem. For the auxiliary construction, whenever a
control meets the hypotheses the same nonvanishing multiplier preserves
its collisions. For the direct construction, a positive success
probability and two valid outcome effects can still have equal signed
probabilities. The exact theta expansion identifies the canonical input;
identifying it is not yet a quantitative separation from a colliding source.

The negative conclusions in this paper are correspondingly limited.
The unbounded positive-Newman multiplier concerns its action on
unrestricted $L^2$. The inverse amplification concerns the fixed free
proper-time response. The absence of a full-cylinder Gibbs state concerns
the retained infinite radial volume. None rules out all enlarged models,
alternative preparations or auxiliary estimates.

\subsection{Quantifiers needed by subsequent analysis}
Three distinctions should accompany every use of the background.
First, all physical angular shells and the full continuous radial
spectrum are retained; cutoffs appear only in justified limiting proofs.
Second, every fixed derivative or moment order is controlled, but its
constant can grow with that order. No order-uniform estimate or infinite
operator expansion is obtained by listing those finite-order statements.
Third, the complex source observation is entire, while the raw scalar
$G$ integral has its initial strip and real bounded measurements have
their real-frequency domain.

The common unresolved mathematical need is a useful signed relation or
estimate with a stated parameter range and a demonstrated advantage over
direct analysis of \eqref{eq:fullH}. Such an estimate is not supplied by
the positive energy moments, the normalized-state tangent, the complete
instrument, or the auxiliary nonvanishing factor alone. The present
manuscript supplies the formulas, domains and exact normalization needed
to formulate and assess that next step.

\appendix
\section{Source and evidence map}\label{sec:evidence}

The following map identifies the proof sources used in this synthesis.
Their full paths, exact input hashes and current review links are recorded
in the manuscript companion README and the TND-E06 input record. Earlier
dated planning paragraphs in the overviews are historical; current
branch dispositions are taken from the latest progress rows and
completed reviews.

\begin{center}\small
\begin{tabularx}{\textwidth}{@{}>{\raggedright\arraybackslash}p{0.24\textwidth}Y@{}}\toprule
Manuscript material & Principal source and evidence carried forward \\\midrule
Geometry and spectra & Parent Note02, Sections 1--3; shared Review01 accepts the stated comparison. Classical angular spectrum and curvature conventions use Camporesi--Higuchi. \\
Green trace, source and tails & Parent Note02, Sections 4 and 6, using MC11's accepted operator/source notes; full theta inversion uses DLMF 20.7.32. \\
Energy dictionary & MC11 observable Note13, Sections 2--4; only its trace/differentiation identities are rederived here. No pending detector or preparation claim is promoted. \\
Scalar field and transport & Parent Note02, Sections 5--7, and the radial companion. Shared Review01 clarifies the two scalar heat generators; the companion's inherited same-assistant evidence remains explicit. \\
Direct construction & Direct Note02, Sections 2--6, and direct Review01. Acceptance is within the supplied ideal-operation scope, from the proof author's same-chat internal assessment. \\
Auxiliary construction & Auxiliary Note01, Sections 2--8, and auxiliary Review01. Acceptance is within the supplied-apparatus scope, from internal same-model assessment in the proof-author chat. \\
Control interpretation & Recorded collision-control cluster and the premise screens in both branch reviews. No new control witness or exclusion is asserted. \\
\bottomrule\end{tabularx}
\end{center}

The source-sum index, physical angular index, reference column, radial
Fourier label and count outcome have different roles. They are never
identified to bypass an infinite sum. Likewise the source energy trace,
the source amplitude, the normalized density and the final signed
observation have different normalizations. The local consistency check
for this manuscript follows those distinctions and is editorial
mathematical checking, not a new independent assessment of the research.
